\section{Introduction}
\label{sec:introduction}

Student's \(t\)-test is one of the most widely used tools for testing
whether a Gaussian mean is zero when the variance is unknown. Its
validity, however, is a fixed-sample guarantee. In many modern
applications the sample size is not fixed in advance: data are inspected
as they arrive, data collection may stop after promising interim
evidence, or sampling may continue after an inconclusive analysis.
Repeated use of ordinary fixed-sample \(p\)-values can then inflate the
Type I error rate. This motivates an anytime-valid analogue of the
one-sample \(t\)-test, whose evidence may be monitored continuously
without invalidating inference.% \citep{Darling1967,Robbins1970,Lai1976,
%Howard2021,Ramdas2023,Grunwald2024}.

At the same time, many scientific problems come with genuine prior
information about plausible effect sizes. Pilot studies, previous
experiments, meta-analyses, or scientific constraints may suggest that
some effects are more plausible than others. Such information can
substantially improve power, especially in sequential analyses where
early rejection is valuable. But informative priors are also risky:
when the prior puts little mass near the direction supported by the
data, a prior-weighted test can be valid but nearly powerless. The goal
of this paper is to construct an anytime-valid \(t\)-test that can use
prior information when it is helpful, while reverting to a robust
noninformative benchmark when the prior is badly wrong.

Let \(Y_1,Y_2,\ldots\) be independent observations from a Gaussian model
with unknown mean \(\mu\) and unknown standard deviation \(\sigma\). We
write $\theta^\star=\mu/\sigma$ for the standardized effect size, and consider
\[
  H_0:\theta^\star=0
  \qquad\text{against}\qquad
  H_1:\theta^\star\neq 0 .
\]
Equivalently, the (composite) null hypothesis is \(\mu=0\), with the nuisance scale
\(\sigma\) left unspecified. We seek a sequence of \(p\)-variables
\((P_n)_{n\geq2}\), or \(p\)-process, satisfying
\[
  \sup_{\sigma>0}
  \mathbb{P}_{\mu=0,\sigma}\!\left(\exists n\geq2: P_n\leq\alpha\right)
  \leq \alpha,\qquad \alpha\in(0,1).
\]
Thus the null may be rejected at the first time \(P_n\leq\alpha\), even
if that time is chosen after looking at the data.

Our construction builds on the work of~\cite{Lai1976,PerezOrtiz2024,Wang2025,Lindon2026} on anytime-valid $t$ tests based on (extended) e-processes. It starts from a nonnegative scale-invariant evidence process \(M_n\) against the composite
null, obtained by mixing a scale-invariant likelihood ratio over the prior distribution on the effect size. In the ordinary case, \(M_n\) is a test martingale under \(H_0\). In the noninformative case it may be an
extended nonnegative martingale arising from a \(\sigma\)-finite mixture~\citep{Lai1976,Wang2025a}. For test martingales,
Ville's inequality \citep{Ville1939} gives the classical
conversion \(P_n=\min\{1,M_n^{-1}\}\). More generally, the extended
Ville inequality~\citep{Lai1976,Wang2025a} yields a decreasing calibration function \(g\), defined
below, such that
\[
  P_n = g(M_n)
\]
is an anytime-valid \(p\)-process.

The evidence processes considered in this paper can also be interpreted as evidence-to-null Bayes factors within a Bayesian \(t\)-test, when using a Jeffreys prior on the unknown scale parameter $\sigma$~\citep{Jeffreys1961,Zellner1980,Rouder2009,Gronau2020}. Such Bayes factors have become a standard Bayesian counterpart of the \(t\)-test in applied work.
% and are implemented in software such as the R package \texttt{BayesFactor}
%\citep{MoreyRouder2015} and JASP \citep{Wagenmakers2018,vanDoorn2021}.
%More general informed Bayesian \(t\)-tests, using flexible priors on
%the standardized effect size, were developed by \citet{Gronau2020}.





%This
%formula is the organising principle of the paper: construct a suitable
%Bayes-factor-like martingale \(M_n\), then calibrate it into a
%sequential \(p\)-value.
%
%The Bayesian \(t\)-test is an especially natural source of such
%evidence processes. The Jeffreys--Zellner--Siow construction
%\citep{Jeffreys1961,Zellner1980} leads to the Cauchy-prior Bayes
%factor popularised by \citet{Rouder2009}, which has become a standard
%Bayesian counterpart of the \(t\)-test in applied work and is
%implemented in software such as the R package \texttt{BayesFactor}
%\citep{MoreyRouder2015} and JASP \citep{Wagenmakers2018,vanDoorn2021}.
%More general informed Bayesian \(t\)-tests, using flexible priors on
%the standardized effect size, were developed by \citet{Gronau2020}.
%From the anytime-valid perspective, these Bayes factors are not merely
%fixed-sample Bayesian summaries. When using a Jeffreys prior on the unknown scale parameter $\sigma$, the corresponding Bayes factor can be interpreted as an inverse test martingale, hence defines a valid p-process.
%
%The unknown-variance case makes this construction nontrivial. One
%cannot simply replace \(\sigma\) by the sample standard deviation in a
%known-variance likelihood ratio. The nuisance scale must be removed by
%invariance or right-Haar integration, and the resulting likelihood
%ratios are martingales only on the appropriate scale-invariant, or
%reduced, filtration. This point is already present in the
%improper-mixture sequential \(t\)-test of \citet{Lai1976}, and has been
%clarified recently through generalized martingales, reduced
%filtrations, extended Ville inequalities, and invariant
%\(e\)-statistics \citep{Wang2025,Wang2025a,PerezOrtiz2024}.

The main contribution of this paper is to determine which priors and calibration function allow
the Bayes-assisted anytime-valid \(t\)-test to be both useful and
robust. The improper Lebesgue mixture recovers the noninformative
anytime-valid \(t\)-test of \citet{Lai1976} and \citet{Wang2025}. For
proper informative priors with polynomial tails, we show that the
extended-Ville \(p\)-process converges back to this noninformative
benchmark as the prior centre is moved arbitrarily far from the data.
Polynomial-tailed priors therefore have a bounded-influence property:
they can improve power under prior--data agreement, yet fall back to
the robust noninformative test under extreme prior--data conflict.
Exponential-tailed priors form a boundary case, while lighter-tailed
priors, such as Gaussian priors, can be fragile under severe
misspecification.

The resulting procedure is Bayesian in how it distributes evidence
across effect sizes, frequentist in its anytime-valid Type I error
guarantee, and robust in a precise tail-dependent sense. It is reminiscent of work in the Bayesian literature on bounded-influence priors \citep{Dawid1973,Pericchi1992,Pericchi1995}. It extends the
bounded/vanishing-influence, extended-Ville programme of \citet{Cortinovis2025}
from the known-variance Gaussian mean problem to the classical and
practically more important unknown-variance \(t\)-testing problem.

%
%
%
%Student's \(t\)-test is one of the most widely used tools for testing
%whether a Gaussian mean is zero when the variance is unknown. Its
%validity, however, is a fixed-sample guarantee. In many modern
%applications the sample size is not fixed in advance: data are inspected
%as they arrive, data collection may stop after promising interim
%evidence, or sampling may continue after an inconclusive analysis.
%Repeated use of ordinary fixed-sample \(p\)-values can then inflate the
%Type I error rate. This motivates an anytime-valid analogue of the
%one-sample \(t\)-test, whose evidence may be monitored continuously
%without invalidating inference.
%
%At the same time, many scientific problems come with genuine prior
%information about plausible effect sizes. Pilot studies, previous
%experiments, meta-analyses, or scientific constraints may suggest that
%some effects are more plausible than others. Such information can
%substantially improve power, especially in sequential analyses where
%early rejection is valuable. But informative priors are also risky:
%when the prior puts little mass near the direction supported by the
%data, a prior-weighted test can be valid but nearly powerless. The goal
%of this paper is to construct an anytime-valid version of the \(t\)-test that can use
%prior information when it is helpful, while reverting to a robust
%noninformative benchmark when the prior is badly wrong.
%
%Let \(Y_1,Y_2,\ldots\) be independent observations from a Gaussian model
%with unknown mean \(\mu\) and unknown standard deviation \(\sigma\). We
%write
%\[
%  \theta^\star=\mu/\sigma
%\]
%for the standardized effect size, and consider
%\[
%  H_0:\theta^\star=0
%  \qquad\text{against}\qquad
%  H_1:\theta^\star\neq 0 .
%\]
%Equivalently, the null hypothesis is \(\mu=0\), with the nuisance scale
%\(\sigma\) left unspecified. We seek a sequence of \(p\)-variables
%\((P_n)_{n\geq2}\), or \(p\)-process, satisfying
%\[
%  \sup_{\sigma>0}
%  \mathbb{P}_{\mu=0,\sigma}\!\left(\exists n\geq2: P_n\leq\alpha\right)
%  \leq \alpha,\qquad \alpha\in(0,1).
%\]
%Thus the null may be rejected at the first time \(P_n\leq\alpha\), even
%if that time is chosen after looking at the data.
%
%Our construction starts from an evidence process \(M_n\) against the
%null. In the ordinary case, \(M_n\) is a nonnegative test martingale
%under \(H_0\); in the more general case needed for noninformative
%mixtures, it is an extended test martingale. With the orientation used
%here, \(M_n\) is also an alternative-versus-null Bayes factor, or the
%inverse of the null-versus-alternative Bayes factor often reported in
%Bayesian testing. Ville's inequality gives the classical conversion
%\(P_n=\min\{1,M_n^{-1}\}\). More generally, the extended Ville
%inequality yields a decreasing calibration function \(g\), defined
%below, such that
%\[
%  P_n = g(M_n)
%\]
%is an anytime-valid \(p\)-process. Equivalently, one may report the
%running version \(g(\max_{2\leq k\leq n}M_k)\). This simple formula is
%the organizing principle of the paper: construct a suitable
%Bayes-factor-like martingale \(M_n\), then calibrate it into a
%sequential \(p\)-value.
%
%The unknown-variance case makes this construction nontrivial. One
%cannot simply replace \(\sigma\) by the sample standard deviation in a
%fixed-\(\sigma\) likelihood ratio. The nuisance scale must be removed
%by invariance or right-Haar integration, and the resulting likelihood
%ratios are martingales only on the appropriate scale-invariant, or
%reduced, filtration. In this reduced experiment, Bayesian \(t\)-test
%Bayes factors obtained by putting a prior on the standardized effect
%size \(\theta^\star\) become mixture test martingales of the form
%\[
%  M_n^\Pi=\int M_n^\theta\,\Pi(\mathrm{d}\theta),
%\]
%where \(\Pi\) is a working prior over effect sizes. The prior therefore
%allocates sequential testing power across alternatives.
%
%The main contribution of this paper is to determine which priors allow
%this allocation to be both useful and robust. The improper Lebesgue
%mixture recovers the noninformative anytime-valid \(t\)-test of
%\citet{Lai1976} and \citet{Wang2025}. For proper informative priors
%with polynomial tails, we show that the extended-Ville \(p\)-process
%converges back to this noninformative benchmark as the prior centre is
%moved arbitrarily far from the data. Polynomial-tailed priors therefore
%have a bounded-influence property: they can improve power under
%prior--data agreement, yet fall back to the robust noninformative test
%under extreme prior--data conflict. Exponential-tailed priors form a
%boundary case, while lighter-tailed priors, such as Gaussian priors,
%can be fragile under severe misspecification.
%
%The resulting procedure is Bayesian in how it distributes evidence
%across effect sizes, frequentist in its anytime-valid Type I error
%guarantee, and robust in a precise tail-dependent sense. It extends the
%bounded-influence, extended-Ville programme of \citet{Cortinovis2025}
%from the known-variance Gaussian mean problem to the classical and
%practically more important unknown-variance \(t\)-testing problem.


%
%Student's \(t\)-test is one of the most widely used tools for assessing
%whether a mean differs from zero when the variance is unknown. Its
%appeal comes from its simplicity: it removes the nuisance scale by
%studentization and produces a \(p\)-value whose null distribution does
%not depend on the unknown variance. This fixed-sample guarantee,
%however, is not the guarantee needed in many contemporary applications.
%In clinical trials, online experiments, laboratory studies and
%sequential data analyses more generally, data are often inspected as
%they arrive. Investigators may stop after observing strong evidence,
%continue after an inconclusive interim analysis, or adapt recruitment
%and sampling decisions to the accumulating data. A \(p\)-value that is
%valid only at a pre-specified sample size need not remain valid under
%such optional stopping.
%
%This motivates the search for an anytime-valid analogue of the
%\(t\)-test: a procedure whose evidence can be monitored
%continuously while preserving Type I error control. But validity under
%optional stopping is only one part of the problem. In many scientific
%settings, the analyst also has genuine prior information about the
%plausible size and direction of the effect. Pilot studies,
%meta-analyses, previous experiments, expert knowledge and physical or
%biological constraints may all suggest that some effect sizes are more
%plausible than others. A sequential test that ignores such information
%can be unnecessarily inefficient, especially when data are costly or
%slow to collect. On the other hand, a procedure that relies too heavily
%on an informative prior can lose most of its power when the prior is
%wrong. The central question of this paper is therefore how to use prior
%information to gain power, while retaining both anytime-validity and
%robustness to severe prior misspecification.
%
%Let \(Y_1,Y_2,\ldots\) be independent observations from a Gaussian
%model with unknown mean \(\mu\) and unknown standard deviation
%\(\sigma\). We write $\theta^\star=\mu/\sigma$ for the standardized effect size, and consider the one-sample testing
%problem
%\[
%  H_0:\theta^\star=0
%  \qquad\text{against}\qquad
%  H_1:\theta^\star\neq 0 .
%\]
%Equivalently, the null hypothesis is \(\mu=0\), with the nuisance scale
%\(\sigma\) left unspecified. The parameter \(\theta^\star\), rather
%than \(\mu\) itself, is the natural target of prior information: it is
%dimensionless, invariant to changes of measurement units, and is the
%quantity on which Bayesian \(t\)-tests and many effect-size
%interpretations are based.
%
%We seek a sequence of \(p\)-variables \((P_n)_{n\geq2}\) satisfying the
%anytime-valid guarantee
%\[
%  \sup_{\sigma>0}
%  \mathbb{P}_{\mu=0,\sigma}\!\left(\exists n\geq2: P_n\leq\alpha\right)
%  \leq \alpha,\qquad \alpha\in(0,1).
%\]
%Such a sequence will be called a \(p\)-process. It can be monitored
%continuously, and the null may be rejected at the first time
%\(P_n\leq\alpha\), without compromising Type I error control. Our goal
%is to construct such \(p\)-processes in a way that also uses prior
%information about \(\theta^\star\). The desired behaviour is
%asymmetric: when the prior is approximately right, the test should
%reject earlier than a noninformative procedure; when the prior is badly
%wrong, the procedure should fall back to a robust noninformative
%benchmark rather than becoming powerless.
%
%The difficulty is that these desiderata pull in different directions.
%Bayesian \(t\)-tests provide a natural way to encode prior information
%about standardized effect sizes, but fixed-sample Bayes factors do not
%by themselves explain how to obtain anytime-valid \(p\)-values under
%unknown variance. Martingale methods provide anytime-valid tests, but
%standard mixture martingales can be highly sensitive to a poorly chosen
%mixing distribution. Noninformative sequential \(t\)-tests, beginning
%with the improper-mixture construction of \citet{Lai1976}, avoid this
%sensitivity but do not exploit genuine prior information. This paper
%connects these strands. We show how informative Bayesian mixtures over
%\(\theta^\star\) can be converted into anytime-valid \(p\)-processes
%for the unknown-variance \(t\)-problem, and we identify exactly which
%prior tails allow the resulting procedure to revert to the
%noninformative test under extreme prior--data conflict.
%
%
%
%The basic martingale idea is simple. A nonnegative test martingale
%\((M_n)\) is a process that starts at one and cannot systematically
%increase under the null. Ville's inequality \citep{Ville1939} implies
%that
%\[
%  P_n^{\mathrm V}=\min\{1,M_n^{-1}\}
%\]
%is a \(p\)-process. Thus a test martingale is a time-uniform measure of evidence against
%\(H_0\): if it ever becomes large, the null can be rejected, regardless
%of when that happens.
%
%The extended Ville inequality gives a sharper and more flexible
%conversion. It applies not only to ordinary integrable test martingales,
%but also to extended nonnegative supermartingales and to
%\(\sigma\)-finite mixtures, including the improper mixtures that arise
%in noninformative sequential testing \citep{Robbins1970,Lai1976,
%Wang2025a}. In the form used here, it provides a decreasing calibration
%function \(g\), defined precisely below, such that the calibrated values
%\[
%  P_n^{eV}=g(M_n)
%\]
%form a valid \(p\)-process. Classical Ville calibration corresponds to
%the crude bound \(g(c)=\min\{1,1/c\}\). The extended calibration can be
%less conservative and, crucially for this paper, remains meaningful for
%the improper mixtures that generate the noninformative anytime-valid
%\(t\)-test.
%
%This martingale viewpoint underlies much of the modern literature on
%confidence sequences, \(e\)-processes and safe anytime-valid inference
%\citep{Darling1967,Robbins1970,Lai1976,Howard2021,Ramdas2023,
%Grunwald2024}. For Gaussian means with unknown variance,
%\citet{Lai1976} gave an early anytime-valid analogue of the \(t\)-test,
%based on improper mixtures. Recent work has clarified this construction
%using extended nonnegative martingales, reduced filtrations and
%extended forms of Ville's inequality \citep{Wang2025,Wang2025a}. In
%parallel, \citet{PerezOrtiz2024} developed the role of group
%invariance and right-Haar integration in producing valid and optimal
%\(e\)-statistics.
%%These results show why the unknown-variance case is
%%not obtained by simply replacing \(\sigma\) by the sample standard
%%deviation in a fixed-\(\sigma\) likelihood ratio. The nuisance scale
%%must first be removed, either by invariance or by integration. The
%%resulting likelihood ratio is a valid martingale only on the
%%appropriate scale-invariant, or reduced, filtration.
%
%The same construction also connects naturally to Bayesian \(t\)-tests.
%Bayesian tests for a Gaussian mean with unknown variance typically put
%a prior on the standardized effect size \(\theta^\star\), together with
%the right-Haar prior \(d\sigma/\sigma\) for the nuisance scale. The
%Jeffreys--Zellner--Siow tradition \citep{Jeffreys1961,Zellner1980}
%leads to the Cauchy-prior Bayes factor popularised by
%\citet{Rouder2009}, while \citet{Gronau2020} develop more general
%informed Bayesian \(t\)-tests based on flexible priors for
%\(\theta^\star\). From the anytime-valid perspective, these Bayes
%factors are not merely fixed-sample Bayesian summaries. With the
%likelihood ratio oriented as alternative versus null and adapted to the
%reduced filtration, they are mixture test martingales of the form
%\[
%  M_n^\Pi=\int M_n^\theta\,\Pi(\mathrm{d}\theta),
%\]
%where \(\Pi\) is a working prior over effect sizes. The prior therefore
%has a direct sequential interpretation: it allocates testing power
%across possible alternatives.
%
%This allocation creates an important robustness question. If \(\Pi\)
%puts substantial mass near the true effect size, the mixture martingale
%can grow quickly and the test can reject earlier than a diffuse or
%noninformative procedure. If \(\Pi\) is badly misspecified, however, the
%same mixture may put little mass in the direction supported by the
%data. The resulting procedure may then be valid but nearly powerless.
%For sequential testing this distinction is essential: validity under
%optional stopping is necessary, but it is not enough. A useful
%Bayes-assisted procedure should exploit reliable prior information
%without collapsing when the prior points to the wrong part of the
%parameter space.
%
%In the known-variance Gaussian mean problem, \citet{Cortinovis2025}
%showed that extended Ville calibration gives precisely such a
%bounded-influence use of informative priors. Their procedures can
%improve efficiency when the prior is well specified, while sufficiently
%heavy-tailed priors protect against arbitrary prior--data conflict. The
%present paper develops the corresponding theory for the
%unknown-variance one-sample \(t\)-test. This is not a formal extension.
%When \(\sigma\) is unknown, the identifiable effect is the ratio
%\(\theta^\star=\mu/\sigma\), the relevant maximal invariant is a
%bounded transformation of the data closely related to the Student
%\(t\)-statistic, and the Bayes factors obtained by mixing over
%\(\theta^\star\) are martingales only after reduction to the
%scale-invariant filtration.
%
%Our first contribution is to construct Bayes-assisted \(p\)-processes
%for the unknown-variance \(t\)-problem. We mix over standardized effect
%sizes, using either proper informative priors or the improper Lebesgue
%mixture, and then convert the resulting ordinary or extended test
%martingales into \(p\)-processes using Ville or extended Ville
%calibration. The improper Lebesgue mixture recovers the
%noninformative anytime-valid \(t\)-test of \citet{Lai1976} and
%\citet{Wang2025}. Proper priors, by contrast, allow the procedure to
%concentrate power near scientifically plausible alternatives.
%
%Our main theoretical contribution is a robustness classification
%governed by the tail behaviour of the effect-size prior. For proper
%informative priors with polynomial tails, we show that, for every fixed
%sample size and dataset, the extended-Ville \(p\)-variable converges to
%the noninformative Lebesgue-mixture benchmark as the prior centre is
%moved arbitrarily far from the data. Polynomial-tailed priors therefore
%have the desired fallback property: they can improve power when the
%prior and the data agree, yet automatically revert to the robust
%noninformative \(t\)-test under extreme prior--data conflict.
%Exponential-tailed priors form a boundary case. For Laplace priors, the
%limiting procedure remains bounded but retains a directional component
%from the prior rather than fully reverting to the Lebesgue benchmark.
%Lighter-tailed priors, represented by Gaussian priors, are fragile
%under severe misspecification and can become effectively powerless.
%
%The resulting test is Bayesian in how it allocates evidence across
%effect sizes, frequentist in its anytime-valid Type I error guarantee,
%and robust in a precise bounded-influence sense. It extends the
%extended-Ville programme of \citet{Cortinovis2025} from the
%known-variance Gaussian mean problem to the classical and practically
%more important unknown-variance \(t\)-testing problem. In doing so, it
%provides a sequential \(t\)-test that can use reliable prior information
%without making its performance depend catastrophically on that
%information being exactly right. 